\documentclass[11pt,a4paper]{article}

\usepackage[T1]{fontenc}
\usepackage[utf8]{inputenc}
\usepackage{lmodern}
\usepackage[margin=25mm]{geometry}
\usepackage{amsmath}
\usepackage{booktabs}
\usepackage{tabularx}
\usepackage{flafter}
\usepackage{listings}
\lstset{basicstyle=\ttfamily\footnotesize,breaklines=true,
    columns=fullflexible,keepspaces=true,showstringspaces=false,
    frame=single,numbers=left,numberstyle=\tiny,tabsize=4}
\usepackage[hidelinks]{hyperref}
\usepackage{microtype}

\newcommand{\localstatrevision}{7f6dec6537cc9e9ca9f9c37961b8b6811c5b7aff}
\hypersetup{
    pdftitle={LocalStat: A Native macOS Environment for Local Execution of a SAS-Language Subset},
    pdfauthor={Muhammad Hafiz bin Ismail},
    pdfsubject={System paper with code and controlled before/after memory and PRINT optimizations}
}

\title{LocalStat: A Native macOS Environment for Local
Execution of a SAS-Language Subset}
\author{Muhammad Hafiz bin Ismail}
\date{4 November 2026\\Draft manuscript with controlled optimization results}

\begin{document}
\maketitle

\begin{abstract}
LocalStat combines a native macOS editing environment with a
local interpreter for a supported subset of the SAS language.
The application uses Swift and AppKit for document editing,
execution control, diagnostics, and result presentation. Its
runtime is a pinned and locally extended version of OpenSAS,
implemented in Zig. Local extensions include a bounded ordinary
least-squares implementation of PROC REG and structured table
output for selected procedures. This paper describes the system
architecture, execution model, numerical implementation, and
compatibility boundaries. An initial 855-sample study is followed
by a controlled 1,912-sample before/after study on an Apple M5
Mac with 16 GB of memory. The implementation changes reclaim
CSV parsing intermediates, eliminate duplicate snapshot row
allocations, and reuse PRINT layout measurements. Across seven
one-million-row workload families, median peak interpreter RSS
falls by 54.0 to 75.6\%, with runtime speedups of
1.12 to 1.26 times. For 10,000 PRINT rows, median table construction
falls from 210.5 to 51.8 ms, and full application latency falls
from 343.5 to 172.9 ms. Outputs and pagination remain unchanged
under the stated checks. The paper includes evaluated code
excerpts and reproducible artifacts. These results apply to the
tested workloads and do not establish general SAS equivalence
or an advantage over another implementation language.
\end{abstract}

\section{Introduction}

LocalStat provides a macOS environment for editing and executing
supported SAS programs locally. The application combines source
editing, project navigation, execution controls, diagnostics,
text output, formatted tables, and report export. Once built,
the application runs without a developer toolchain or a hosted
execution service.

The system uses OpenSAS as its interpreter foundation~\cite{opensas}.
LocalStat adds a native application layer and
selected runtime extensions. The implementation is independent
of SAS Institute's runtime and supports a bounded language
subset. Compatibility must therefore be described through
specific supported behavior rather than inferred from the
ability to execute a program successfully.

This paper documents three implementation areas: the integration
of a native editor with a local interpreter, a bounded regression
procedure, and a structured interface for presenting selected
procedure results. It reports a reproducible baseline followed
by targeted memory and PRINT optimizations with a controlled
before/after evaluation.

\section{System Architecture}

\subsection{Native application}

The application is implemented in Swift using AppKit. Its
interface includes document tabs, a project sidebar, an editor,
and results views. Users can run an entire editor buffer,
including unsaved changes, with the selected project directory
serving as the runtime working directory.

The application retains the executed source snapshot with the
run result. Diagnostic navigation is enabled only when an
original-source position refers to an unchanged editor buffer.
Positions belonging to macro-expanded text are not treated as
direct coordinates in the original source.

Offline Help documents supported behavior and provides executable
examples. Text export preserves runtime output and diagnostics.
PDF report export combines the executed source, run metadata,
results, and diagnostics. This report facility belongs to the
application and is distinct from a SAS ODS destination.

\subsection{Interpreter process}

Each submission starts a fresh interpreter process. The
application writes the source snapshot to a temporary file,
launches the bundled executable, and drains standard output and
standard error concurrently. Concurrent draining prevents a full
output pipe from blocking completion.

Fresh processes reset WORK datasets and macro definitions between
submissions. Files explicitly written by a program can persist.
The process boundary separates run state; it does not provide a
filesystem sandbox.

The application supports cancellation through a termination
signal, followed by forced termination after two seconds if the
process remains active. Runtime output is buffered until
execution finishes, so cancellation can discard output that has
not yet been emitted.

\subsection{Runtime provenance and build}

The interpreter is based on OpenSAS commit
\texttt{927a2fa9f1401e365d975536bd704819ada1f1f4}.
The repository retains the upstream source snapshot and records
local modifications separately~\cite{runtimechanges}.

The current build targets Apple Silicon and macOS 14 or newer.
The runtime is built with Zig 0.16.0 using
\texttt{ReleaseSafe}; the Swift application uses a release build.
Python supports repository tests and Help generation. It is not
part of the application execution path.

\section{Selected Implementation Details}

\subsection{Bounded ordinary least squares}

LocalStat extends the runtime with a numeric subset of
\texttt{PROC REG}. A supported model has one response, explicit
numeric predictors, and an intercept:
\[
    y = X\beta + \varepsilon.
\]
Rows missing any model variable are excluded. The remaining
design must have positive residual degrees of freedom and
sufficient numerical rank.

The solver uses column scaling and Householder QR factorization
with column pivoting. It produces coefficient estimates,
standard errors, analysis-of-variance quantities, and associated
test statistics.

The implementation rejects numerically rank-deficient designs.
It does not reproduce generalized-inverse handling of aliased
coefficients. Unsupported model options produce explicit
diagnostics. This procedure is a documented subset rather than
a complete implementation of SAS PROC REG~\cite{compatibility}.

\subsection{Structured table results}

The runtime can write an optional, versioned JSON sidecar for
selected results from \texttt{PROC PRINT},
\texttt{PROC FREQ}, and \texttt{PROC TABULATE}. The sidecar
contains table titles, headers, rows, notes, and cell spans.

The application validates the version and table shape before
displaying the results. Raw output and diagnostics remain
available alongside formatted tables. The interface supports
native table presentation without interpreting aligned text as
a table schema.

\section{Evaluation Questions}

The evaluation addresses four questions:
\begin{enumerate}
    \item How do interpreter runtime and peak memory change as
    input size increases?
    \item Which runtime functions and execution stages account
    for the largest costs?
    \item How much additional latency comes from table generation,
    result decoding, and application presentation?
    \item Do targeted optimizations improve performance while
    preserving validated behavior?
\end{enumerate}

The study does not attempt to establish that one implementation
language is universally faster than another. Such a comparison
would require equivalent implementations and controlled
workloads. The first three questions are addressed through
baseline measurements and profiles. Section~\ref{sec:optimization}
addresses the fourth with a controlled implementation comparison.

\section{Experimental Method}

\subsection{Workloads}

A six-observation multiple-regression example from the repository,
with an appended checksum, provides a small-program baseline.
Synthetic workloads use 1,000, 10,000, 100,000, and 1,000,000 rows.
Every generated CSV has five numeric columns: a descending row
identifier, a group identifier, two predictors, and a response.
The group identifier cycles through 100 levels. Predictors and
noise cycle through a balanced factorial design with values
$-1$ and $1$, and the response is
$y=2+3x-4z+\varepsilon$.

The seven scaling families perform a DATA-step transformation
($v=3\,\mathrm{id}+2$), ascending sort, SQL grouping, SQL join,
two-predictor regression, CSV round trip, or frequency tabulation.
The grouping workload produces 100 groups. The join uses a
100-row dimension with unique group keys and produces exactly
as many rows as its input. The CSV round trip imports, exports,
and reimports all five columns. PRINT adds output-heavy cases
at 1,000 and 10,000 rows. Both FREQ and PRINT run with and
without the structured table sidecar.

Input generation occurs outside the timed interval. Each scaling
program includes CSV loading and a compact validation scan;
these measurements are whole-program costs, not isolated
operator timings. Checks verify row counts and exact checksums,
sort order, group aggregates, and joined key associations.
Regression checks require coefficients $(2,3,-4)$ within absolute
and relative tolerances of $10^{-8}$, and residual sum of squares
equal to the row count within relative tolerance $10^{-8}$.
Every exported CSV field and numeric listing row is checked
after timing. Sidecar checks validate shape, PRINT cells, and
FREQ group counts. Output and diagnostic hashes must remain identical
across repetitions and sidecar variants.

\subsection{Measurement boundaries}

Command-line measurements run the executable directly. A
high-resolution monotonic clock surrounds process launch and
\texttt{wait4}; elapsed time includes launch and observation
overhead, source loading, interpretation, and output completion.
Per-child \texttt{wait4} resource counters provide CPU time and
peak resident memory. This replaces the draft protocol's
\texttt{/usr/bin/time -l} invocation to retain sub-centisecond
precision and avoid a separate timing process. Compilation,
input generation, output parsing, and validation are excluded.

The native experiment compares the existing
\texttt{LocalRuntimeRunner} directly against the actual
\texttt{AppDelegate.runProgram} workflow in a release-mode XCTest
host. Both execute the same frozen interpreter. The direct
runner includes source snapshot creation, process handling,
output capture, and table decoding. Application timing starts
before dispatch and ends after result application, layout, and
a synchronous window display pass. Completion observation polls
at one-millisecond intervals. Window creation and document setup
are excluded. The four cases are DATA at 10,000 rows, FREQ at
100,000 rows, and PRINT at 1,000 and 10,000 rows.

These are paired executions, not separately instrumented stages
of one execution. Their difference includes scheduling and
observation effects and cannot isolate decoding or rendering
time. The display boundary does not measure physical screen
presentation or the latency of a human keyboard action.

Separate one-second stack profiles were collected at nominal
one-millisecond sampling intervals for each million-row family.
Profiled runs are excluded from timing summaries. Their short,
partially observed execution windows support hotspot
identification rather than precise stage-time attribution.

\subsection{Repetition and environment}

Each configuration has one separately recorded first execution,
three warmups, and fifteen measured repetitions in fresh
interpreter processes. The primary summaries are medians and
interquartile ranges (IQRs), using linearly interpolated inclusive
quartiles. The first runs and warmups are retained in the raw
artifact but excluded from summary tables. No cache flush was
performed, and first execution is not claimed to be a cold-cache
measurement. Paired variant order alternates by iteration.

The host was a ten-core Apple M5 MacBook Pro (Mac17,2) with 16 GB
of memory, running macOS 26.4.1 (25E253), Swift 6.3.3, and Zig
0.16.0. The Mac was connected to AC power with low-power mode
disabled. The runtime used the normal ARM64 ReleaseSafe build;
the native test host used Swift release compilation. App windows
used the light appearance, a $1100\times780$ frame, default fonts,
and unit table zoom. Timing studies and profiling ran
sequentially.

The artifact records executable hashes, source hashes, working-tree
state, compiler versions, power settings, and load averages.
Native app sources were frozen before compilation to prevent
concurrent workspace edits from changing that experiment. The
machine was not otherwise isolated from operating-system or user
activity. Command-line runs used the automation filesystem
sandbox; native app runs required normal AppKit access. The
native runner and application comparison used the same host
process and access conditions.

\subsection{Correctness and comparisons}

A sample is accepted only when execution succeeds and all
applicable checks pass. A sixty-second command-line watchdog
kills the process group on timeout; failed and timed-out runs
are recorded rather than included as successful observations.
Native output hashes are checked against validated command-line
output. The reporting tool recomputes summaries from raw records
and rejects incomplete, duplicated, or invalid repetitions.

For paired sidecar and native comparisons, the reported
difference is the median of iteration-matched differences,
not the difference of marginal medians. No significance test
or confidence interval is asserted. An IQR crossing zero is
treated as insufficient evidence of a consistent directional
effect in that comparison. The initial baseline did not compare
optimized builds; the follow-up comparison is described separately.

\section{Initial Baseline Results}

\subsection{Interpreter latency and memory}

All 703 command-line samples passed validation, with no timeouts.
There are 37 configurations and 555 measured samples; the other
148 are first executions or warmups. Table~\ref{tab:cli} reports
the three larger sizes. The six-row example has a median latency
of 2.7 ms (IQR 0.7 ms), including launch and observation costs.

\begin{table}[tbp]
\centering
\small
\begin{tabular}{lrrrr}
\toprule
Workload & 10,000 & 100,000 & 1,000,000 & MiB \\
\midrule
DATA transform & 12.0 (0.7) & 107.6 (6.8) & 1035.0 (52.5) & 1832.7 \\
SORT & 10.7 (4.6) & 85.7 (4.0) & 877.2 (29.1) & 1533.0 \\
SQL grouping & 10.3 (0.9) & 93.2 (5.5) & 883.4 (21.9) & 1614.7 \\
SQL join & 15.5 (0.7) & 137.8 (3.3) & 1346.6 (62.4) & 2005.5 \\
REG (two predictors) & 10.8 (0.7) & 94.1 (4.3) & 1045.0 (100.3) & 1625.4 \\
CSV round trip & 16.7 (2.2) & 152.6 (10.3) & 1677.8 (273.0) & 1994.2 \\
FREQ & 10.0 (0.6) & 80.3 (1.8) & 949.2 (109.6) & 1464.5 \\
\bottomrule
\end{tabular}

\caption{Command-line median latency in milliseconds, with IQR
width in parentheses; fifteen measured repetitions per cell.
The last column is median peak resident memory in MiB at one
million rows. Sidecar output is disabled.}
\label{tab:cli}
\end{table}

At one million rows, median latency ranges from 877.2 ms for SORT
to 1677.8 ms for the CSV round trip. Increasing input from
100,000 to one million rows increases median latency by factors
of approximately 9.5 to 11.8 across these families. These
observations describe this input range; they do not establish
asymptotic complexity.

Median peak resident memory at one million rows ranges from
1464.5 MiB for FREQ to 2005.5 MiB for the SQL join. The shared
CSV load, dataset representation, and validation scans are part
of each workload. The table therefore cannot attribute either
time or memory solely to a named procedure. The CSV round trip
has a 273.0 ms latency IQR at one million rows, illustrating
variation that a single execution would conceal.

\subsection{Structured table output}

Table~\ref{tab:sidecar} compares the sidecar variants. The paired
median increase for PRINT is 0.6 ms at 1,000 rows and 4.3 ms at
10,000 rows. FREQ produces only 100 category rows regardless of
input size. For FREQ at 10,000 through one million input rows,
the paired-difference IQR crosses zero, so these observations
do not establish a consistent directional sidecar cost.

\begin{table}[tbp]
\centering
\small
\begin{tabular}{lrrr}
\toprule
Workload & Plain & Sidecar & Paired $\Delta$ \\
\midrule
FREQ (1,000) & 2.8 (0.1) & 3.0 (0.1) & 0.2 (0.2) \\
FREQ (10,000) & 10.0 (0.6) & 10.0 (0.9) & -0.2 (1.0) \\
FREQ (100,000) & 80.3 (1.8) & 80.7 (1.6) & 0.0 (1.0) \\
FREQ (1,000,000) & 949.2 (109.6) & 932.6 (100.7) & 14.1 (68.9) \\
PRINT (1,000) & 3.6 (0.4) & 4.2 (0.3) & 0.6 (0.3) \\
PRINT (10,000) & 16.1 (3.0) & 20.7 (0.9) & 4.3 (2.0) \\
\bottomrule
\end{tabular}

\caption{CLI table-output comparison in milliseconds. Values
are medians (IQR widths). The paired difference is sidecar minus
plain execution, matched by iteration. It can differ in sign
from the difference of marginal medians.}
\label{tab:sidecar}
\end{table}

\subsection{Native application latency}

All 152 native samples passed checks, including 120 measured
samples across eight configurations. Table~\ref{tab:ui} compares
the direct native runner and complete app workflow. Native runner
costs include Foundation process management and table decoding
and are distinct from the command-line launch path.

\begin{table}[tbp]
\centering
\small
\begin{tabular}{lrrr}
\toprule
Workload & Native runner & Application & Paired $\Delta$ \\
\midrule
DATA transform (10,000) & 65.9 (2.2) & 74.1 (2.3) & 7.7 (3.6) \\
FREQ (100,000) & 134.8 (3.3) & 143.8 (6.9) & 8.0 (5.7) \\
PRINT (1,000) & 71.5 (1.5) & 104.9 (2.7) & 32.9 (5.3) \\
PRINT (10,000) & 107.4 (2.2) & 385.1 (75.7) & 276.4 (76.3) \\
\bottomrule
\end{tabular}

\caption{Native runner and application latency in milliseconds:
medians (IQR widths), fifteen measured repetitions per cell.
The paired difference is application minus runner. App timing
ends after the first synchronous display pass, not physical
screen presentation.}
\label{tab:ui}
\end{table}

For the DATA and FREQ cases, paired median application-minus-runner
differences are 7.7 and 8.0 ms. For PRINT, the difference grows
from 32.9 ms at 1,000 rows to 276.4 ms at 10,000 rows. The latter
case has a median application latency of 385.1 ms and an IQR of
75.7 ms. This identifies result-heavy application behavior as a
separate optimization target, but does not isolate the cost of
formatting, pagination, layout, drawing, or scheduling.

\subsection{Profiles and optimization priorities}

All seven profiled million-row runs completed successfully.
The samples repeatedly show \texttt{io.readDelimited}, numeric
conversion, PDV field lookup, hash lookup, and memory allocation
or initialization. The DATA profile also shows row consumption,
expression evaluation, and dataset output. These observations
support investigating CSV conversion, field access, and data
representation before considering a language migration. They
do not quantify an achievable speedup or establish a precise
breakdown of total execution time.

\section{Targeted Optimizations and Controlled Comparison}
\label{sec:optimization}

\subsection{Allocation lifetime and row ownership}

The CSV reader previously decoded all records into temporary
string fields allocated in the run arena, then constructed
typed rows while retaining those temporary fields. Numeric
columns therefore kept both their textual intermediates and
their final values until the process exited. DATA-step snapshots
also allocated a temporary value row and passed it to
\texttt{Dataset.appendRow}, which allocated and copied another row.

The revised importer makes two passes over its in-memory record
index. The first preserves whole-column numeric and date
inference while resetting a record-local arena after each row.
The second constructs owned values, again reusing record-local
scratch space. Temporary record-index storage is released when
import returns. Character cells are explicitly copied into the
destination allocator before scratch space is reset. This does
not make file loading fully streaming: source text, the record
index, and final datasets still require memory proportional to
input size.

The snapshot path now allocates final cells once, copies character
bytes into dataset-owned storage, and appends that already-owned
row directly. Listing~\ref{lst:snapshot} shows the complete helper.
The remaining listings are excerpts from the exact evaluated
source, preserved with file and line provenance in the artifact.

\noindent\begin{minipage}{\linewidth}
\lstinputlisting[caption={Single-allocation PDV snapshot with owned character bytes.},label={lst:snapshot}]{docs/paper/results/2026-11-04-optimization/code/snapshot.zig}
\end{minipage}

\noindent\begin{minipage}{\linewidth}
\lstinputlisting[gobble=4,caption={CSV inference pass using reusable per-record scratch storage.},label={lst:inference}]{docs/paper/results/2026-11-04-optimization/code/csv-inference.zig}
\end{minipage}

\noindent\begin{minipage}{\linewidth}
\lstinputlisting[gobble=4,caption={Typed row construction copies character data before reusing scratch storage.},label={lst:materialization}]{docs/paper/results/2026-11-04-optimization/code/csv-materialization.zig}
\end{minipage}

Each allocation-budget test constructs 4,096 five-column numeric
rows within a one-MiB destination allocator. Both import and
snapshot tests fail before and pass after the changes. They bound
destination allocation, not process RSS. Mutation tests separately
verify character ownership after input buffers change.

\subsection{PRINT layout reuse}

The original table view repeatedly measured cell text during
pagination and drawing and rebuilt column projections even
when all columns were displayed. The revised view caches text
heights by text, available width, font name, and font size during
construction. It records header and row heights once per column
panel and reuses those arrays for pagination and drawing. Full
column projections reuse the existing row cells. The temporary
text cache is discarded after construction; shared height arrays
remain with the page fragments.

\noindent\begin{minipage}{\linewidth}
\lstinputlisting[gobble=4,caption={Text-height cache keyed by the complete measurement context.},label={lst:heightcache}]{docs/paper/results/2026-11-04-optimization/code/text-height.swift}
\end{minipage}

\noindent\begin{minipage}{\linewidth}
\lstinputlisting[gobble=16,caption={Header and row metrics are computed once per column panel.},label={lst:rowheights}]{docs/paper/results/2026-11-04-optimization/code/row-heights.swift}
\end{minipage}

A deterministic test lays out 2,000 repeated two-cell rows with
at most four distinct text measurements and requires drawing
to perform no additional measurements. Another test distinguishes
font and width contexts. Existing tests cover zoom geometry,
accessibility text, Unicode output, spanning headers, multiple
column panels, and searchable PDF pagination.

\subsection{Follow-up protocol}

The follow-up preserves the initial evidence and uses new frozen
before/after artifacts. Only the table view differs between
the native packages; their measurement driver and remaining
application sources are identical. The baseline uses the source
snapshot taken before these changes. Historical results remain
a separate study. Interpreter changes are confined to \texttt{io.zig}.

For every CLI workload, sidecar variant, and iteration, baseline
and optimized executions are adjacent; their order alternates
between iterations. All four row counts and the six-observation
example are repeated, with one first execution, three warmups,
and fifteen measured repetitions per build and configuration.
Both builds use the same inputs, compiler versions, optimization
settings, and execution access conditions.
The follow-up uses normal local process access for both CLI and
native paths; the initial CLI study used the automation filesystem sandbox.

Native measurements use counterbalanced blocks in the order
before, after, after, before. Each build contributes eight measured
repetitions in its first block and seven in its second. Every
block records a separate first execution and three warmups per
configuration. In addition to the native runner and full app,
FREQ and PRINT measure table-view construction from a preloaded
validated result. This construction measurement excludes runtime
execution, JSON decoding, and drawing. Its page count must be
identical across builds. App timing retains the original
synchronous-display boundary.

The follow-up tables report median latency with IQR width,
median peak child RSS, and ratios of before/after medians.
Speedup is $t_{\mathrm{before}}/t_{\mathrm{after}}$; memory
reduction is $100(1-m_{\mathrm{after}}/m_{\mathrm{before}})$.
No confidence interval or universal complexity claim is inferred.

\subsection{Before/after results}

All 1,406 interleaved CLI observations passed validation, including
1,110 measured observations across 37 configurations per build.
All 506 native observations also passed, including 330 measured
observations across 11 configurations per build. The remaining
observations are first executions and warmups. Output and
diagnostic hashes match across builds and repetitions; the
canonical sidecar JSON hashes also match. Native table page
counts are unchanged: 39 and 385 for the two PRINT sizes, and
four for FREQ.

\begin{table}[tbp]
\centering\small
\begin{tabular}{lrrr}
\toprule
Workload & Before (MiB) & After (MiB) & Reduction (\%) \\
\midrule
DATA transform & 1791.9 & 565.8 & 68.4 \\
SORT & 1533.0 & 403.4 & 73.7 \\
SQL grouping & 1614.7 & 668.2 & 58.6 \\
SQL join & 2005.5 & 921.7 & 54.0 \\
REG & 1625.4 & 518.7 & 68.1 \\
CSV round trip & 2324.0 & 579.7 & 75.1 \\
FREQ & 1464.5 & 357.8 & 75.6 \\
\bottomrule
\end{tabular}

\caption{Million-row median peak child RSS before and after the
allocation changes; fifteen measured repetitions per build and
workload. Percentage reduction compares medians.}
\label{tab:memorycomparison}
\end{table}

Table~\ref{tab:memorycomparison} shows median peak RSS reductions
of 54.0\% for the SQL join through 75.6\% for FREQ. Optimized
medians range from 357.8 to 921.7 MiB, compared with 1464.5 to
2324.0 MiB for the fresh baseline. The CSV round trip falls from
2324.0 to 579.7 MiB. Final datasets and source text still occupy
memory, and SQL and regression retain their own working data.
The changes reduce redundant allocation rather than imposing
a fixed memory cap.

\begin{table}[tbp]
\centering\small
\begin{tabular}{lrrr}
\toprule
Workload & Before & After & Speedup ($\times$) \\
\midrule
DATA transform & 1029.8 (49.0) & 826.6 (8.0) & 1.25 \\
SORT & 827.1 (19.1) & 669.0 (10.5) & 1.24 \\
SQL grouping & 871.6 (36.2) & 750.1 (7.0) & 1.16 \\
SQL join & 1298.2 (22.5) & 1160.8 (7.2) & 1.12 \\
REG & 860.0 (4.9) & 725.6 (8.8) & 1.19 \\
CSV round trip & 1432.5 (41.4) & 1132.7 (12.0) & 1.26 \\
FREQ & 730.2 (4.3) & 598.2 (6.4) & 1.22 \\
\bottomrule
\end{tabular}

\caption{Million-row CLI latency in milliseconds: median (IQR
width). Speedup is the ratio of before/after medians, with
adjacent build order alternating each iteration.}
\label{tab:latencycomparison}
\end{table}

All seven million-row families improve in this controlled run
(Table~\ref{tab:latencycomparison}), with median speedups from
1.12 for the SQL join to 1.26 for the CSV round trip. The
optimized medians range from 598.2 to 1160.8 ms. The importer
parses each record twice, but avoids retaining its intermediate
field buffers; the observed result is lower memory use with
lower elapsed time for these workloads. These measurements do
not isolate allocation, parsing, or memory-page costs individually.

\begin{table}[tbp]
\centering\footnotesize
\begin{tabular}{lrrrr}
\toprule
Workload & Path & Before & After & Speedup ($\times$) \\
\midrule
DATA transform (10,000) & app & 74.2 (1.1) & 74.9 (1.5) & 0.99 \\
DATA transform (10,000) & runner & 69.9 (5.3) & 69.5 (4.3) & 1.00 \\
FREQ (100,000) & app & 142.2 (1.8) & 139.9 (1.2) & 1.02 \\
FREQ (100,000) & layout & 2.6 (0.1) & 1.7 (0.2) & 1.59 \\
FREQ (100,000) & runner & 143.6 (2.6) & 73.4 (33.4) & 1.96 \\
PRINT (1,000) & app & 106.0 (6.6) & 88.1 (7.1) & 1.20 \\
PRINT (1,000) & layout & 22.3 (1.1) & 6.1 (0.5) & 3.67 \\
PRINT (1,000) & runner & 76.4 (4.5) & 76.2 (3.6) & 1.00 \\
PRINT (10,000) & app & 343.5 (41.4) & 172.9 (5.8) & 1.99 \\
PRINT (10,000) & layout & 210.5 (11.3) & 51.8 (1.8) & 4.07 \\
PRINT (10,000) & runner & 114.8 (4.0) & 112.1 (2.6) & 1.02 \\
\bottomrule
\end{tabular}

\caption{Counterbalanced native comparison in milliseconds:
median (IQR width), fifteen measured observations per build and
path. The layout path constructs a table view from preloaded
results and excludes execution, decoding, and drawing.}
\label{tab:nativecomparison}
\end{table}

For PRINT at 10,000 rows, isolated table construction improves
from 210.5 to 51.8 ms, a 4.07-fold speedup
(Table~\ref{tab:nativecomparison}). Full application latency
improves from 343.5 to 172.9 ms, a 1.99-fold speedup. At 1,000
PRINT rows, construction improves from 22.3 to 6.1 ms and the
application improves from 106.0 to 88.1 ms. These comparisons
retain every table row and the original page geometry.

Not every native path shows a meaningful improvement. DATA app
latency changes from 74.2 to 74.9 ms, within the observed IQRs,
and FREQ app latency changes only from 142.2 to 139.9 ms. The
optimized FREQ runner has an IQR of 33.4 ms, so its marginal
median is especially sensitive to scheduling and host behavior.
Runner and app measurements are separate executions; differences
between them must not be interpreted as exact stage costs.

The four source changes discussed here were evaluated together
as one candidate: reusable CSV record storage, single-allocation
snapshots, cached text heights, and stored row/panel metrics.
The study does not provide an ablation assigning separate speedups
to each change. Exact before/after sources, the patch, raw
measurements, and code-excerpt provenance are preserved in the
optimization artifact~\cite{optimization}.

\section{Verification and Reproducibility}

Repository documentation remains pinned to base source revision
\texttt{7f6dec6} (full identifier:
\texttt{\localstatrevision}). The evaluated working tree also
contained uncommitted app changes. The artifact preserves exact
native sources, per-file hashes, runtime executable hashes,
workload source, input hashes, individual samples, and profiles
rather than claiming a clean evaluated release commit.

For the initial study, the project verification log records 969 runtime unit
tests, 1,924 corpus programs, and 11 end-to-end fixtures passing.
The Swift suite reports 48 tests with one opt-in benchmark test
skipped and no failures; the benchmark test was subsequently
enabled for the native study. The packaged-runtime checks,
harness checks, Help checks, and repository text check also
passed. Separate evidence-auditor tests check incomplete records,
output mismatches, and iteration pairing. Earlier verification
records remain historical evidence~\cite{verification}.

For the optimization follow-up, the full project suite passes
973 runtime unit tests, all 1,924 corpus programs, and all 11
end-to-end fixtures. The Swift suite reports 57 tests with the
one opt-in measurement test skipped and no failures. Both
frozen native packages then pass the enabled measurement test.
The packaged-runtime checks and the Python harness/report tests
also pass. The source archive and patch identify the evaluated
working-tree changes without inventing a release commit.

The measurement artifact is stored under
\path{docs/paper/results/2026-11-04}. The harness generates inputs
and runs the CLI study, an opt-in native test measures the app,
and a reporting script audits raw evidence and generates the
three tables in this paper. Instructions describe rebuilding and
repeating the experiments. These checks establish reproducibility
of the stated workflow, not correctness for every SAS program.

\clearpage
\section{Limitations}

LocalStat supports a subset of the SAS language and procedure
options. Passing repository tests does not establish general
SAS compatibility. Existing numerical checks include analytical
and published-example comparisons; this manuscript reports no
new comparison against an installed SAS engine.

Fresh processes provide no persistent interactive session.
Buffered output and in-memory results affect cancellation and
large-output workloads.

The alpha is not sandboxed: programs can access permitted files
and inherited environment values. Local execution alone does not
establish security or suitability for regulated work.

One host, numeric inputs, 100 groups, a bounded join, and two
predictors do not represent all programs. Character-heavy files
and tables with mostly distinct text may behave differently.
The transient text cache can grow during construction; native
process peak memory was not measured.

Warm caches, workload order, OS activity, and execution access
conditions affect results. Native timings use a release XCTest
host and a synchronous display boundary, not independently
launched app or physical-screen latency. Short profiles miss
parts of execution. There is no persistent-session, alternative-
language, installed-SAS, or per-change ablation comparison.
Fresh controls account for differences from historical host and
allocator conditions without assuming old and new baselines match.

\section{Conclusion}

LocalStat integrates a native macOS editing environment with
a locally extended OpenSAS interpreter. Its current design
provides fresh-process execution, source-aware diagnostics,
bounded ordinary least squares, and structured presentation
for selected tables.

The initial baseline identified interpreter memory and large
PRINT presentation costs. The controlled follow-up reduces
million-row median peak interpreter RSS by 54.0 to 75.6\% and
improves median runtime by factors of 1.12 to 1.26 in the tested
families. Cached PRINT layout provides a 4.07-fold construction
speedup at 10,000 rows and a 1.99-fold improvement in full app
latency, with matching output and page counts. The implementation,
regression tests, and preserved artifacts make these specific
results inspectable. They establish a targeted improvement on
this host, not general SAS equivalence or a language advantage.

\begin{thebibliography}{9}
\raggedright

\bibitem{opensas}
OpenSAS contributors.
\emph{OpenSAS source code}.
Pinned commit:
\texttt{927a2fa9f1401e365d975536bd704819ada1f1f4}.
\href{https://github.com/kirha-ai/opensas/tree/927a2fa9f1401e365d975536bd704819ada1f1f4}{OpenSAS source snapshot}.

\bibitem{runtimechanges}
LocalStat project.
\emph{Runtime provenance and local changes}.
\href{https://github.com/hafiz122/localstat/blob/\localstatrevision/docs/RUNTIME_CHANGES.md}{\texttt{docs/RUNTIME\_CHANGES.md}}.
Source revision \texttt{7f6dec6}.

\bibitem{compatibility}
LocalStat project.
\emph{LocalStat source-build compatibility}.
\href{https://github.com/hafiz122/localstat/blob/\localstatrevision/docs/COMPATIBILITY.md}{\texttt{docs/COMPATIBILITY.md}}.
Source revision \texttt{7f6dec6}.

\bibitem{verification}
LocalStat project.
\emph{Verification record}.
\href{https://github.com/hafiz122/localstat/blob/\localstatrevision/docs/VERIFICATION.md}{\texttt{docs/VERIFICATION.md}}.
Source revision \texttt{7f6dec6}; historical records dated
29 October and 1 November 2026.

\bibitem{optimization}
Muhammad Hafiz bin Ismail.
\emph{LocalStat controlled optimization artifact}.
4 November 2026.
\path{docs/paper/results/2026-11-04-optimization}.
Frozen sources, implementation patch, code provenance, individual
observations, generated comparisons, and verification records.

\end{thebibliography}
\end{document}
